\navid{System} verification tools are crucial for ensuring correctness prior to testing, implementation, or deployment, particularly in expensive, high-risk, or safety-critical systems \cite{zhang2023reachability,schilling2022verification,komendera2012intelligent}. In real-world implementations, we face two significant challenges for \navid{system} verification: (1) assuming access to an underlying model for the system is \navid{often prohibitive}, and (2) the presence of noise and uncertainty results in \navid{stochastic systems}, requiring us to perform \navidd{statistical} verification \navidd{to obtain} probabilistic \navidd{safety certificates}. Data-driven statistical \navidd{reachability analysis} is a well-established tool for \navidd{statistical} verification, \navidd{and has been applied} to safety-critical problems across various domains, \navidd{including, autonomy, medical imaging and CPS. } \cite{devonport2021data,fisac2018general,dryvr,hashemi2024statistical}. 

In the context of data-driven statistical reachability analysis, given a user-specified threshold $\delta \in (0,1)$, the goal is to produce a set that \navidd{ensures that} any trajectory during deployment lies within \navidd{this set} with a probability of \navidd{no less than} $\delta$. \navidd{As explained in the "\textit{Related Work}" section, our work is motivated by the method proposed in \cite{hashemi2024statistical}. This methodology involves three }main steps: (1) learning a deterministic surrogate model from sampled trajectories, (2) performing reachability analysis on the surrogate model, and (3) using robust conformal inference  \cite{cauchois2020robust} to calculate an inflating hypercube. This inflating hypercube quantifies the necessary expansion of the surrogate model's reachable set to provide provable probabilistic guarantees on the flowpipe that is obtained for the black-box oracle. 
% The main motivation for the choice of $\relu$ activation function is the recent advancements in reachability analysis of neural networks with ReLU activations. In this case the NNV toolbox from \cite{tran2020nnv} can be employed to compute the reachable set of the surrogate model, utilizing two main methods: (1) the exact-star method, which offers precise but slow computations, and (2) the approx-star method, which is more efficient but can be a source of conservatism.


% Reachability analysis is a popular method to give safety guarantees for
% stochastic cyber-physical systems (SCPSs) that takes in a symbolic description of the
% system dynamics and uses set-propagation methods to compute an overapproximation
% of the set of reachable states over a bounded time horizon.

%  In this paper, we
% investigate the problem of performing reachability analysis for an SCPS that
% does not have a symbolic description of the dynamics, but instead is described
% using a digital twin model that can be simulated to generate system trajectories.
% An important challenge is that the simulator implicitly models a probability 
% distribution over the  set of trajectories of the SCPS; however, it is typical 
% to have a sim2real gap, i.e., the actual distribution of the trajectories in a
% deployment setting may be shifted from the distribution assumed by the simulator.

% The authors in \cite{hashemi2023data, hashemi2024statistical} propose a statistical reachability analysis technique that, given a
% user-provided threshold $\delta \in (0,1)$, provides a set that guarantees that any reachable state during deployment lies in this set with probability not smaller than this threshold. This method is based on three main steps: (1) learning a deterministic surrogate model from sampled
% trajectories, (2) conducting reachability analysis over the surrogate model, and (3) employing {\em robust conformal inference} using an additional set of sampled trajectories to compute for an inflating hypercube that quantifies the extent to which the reachable set of the surrogate model should be inflated to provide provable probabilistic guarantees on the resulting flowpipe. Due to the recent achievements in reachability analysis on neural networks with ReLU activations, they strictly recommend using $\relu$ neural networks as the surrogate models, and rely on the NNV toolbox from \cite{tran2020nnv} to compute the reachable set of the surrogate model. This approach employs two main methods for reachability: (1) the exact-star method, which performs precise but slow computations, and (2) the approx-star method, which is more efficient but more conservative.

% In this paper, we improve the scalability and the accuracy of the proposed approach for data driven reachability in \citep{hashemi2023data,hashemi2024statistical}, that is still suffering from lack of accuracy and scalability over longer time horizon and higher dimensional systems.

% The authors in \cite{hashemi2023data} suggest training the simulated trajectories $\trajsim_{\statee_0} \in \mathbb{R}^{n(\horizon+1)}$ with the $\relu$ NN surrogate model $\overallf(\statee_0 ; \theta)$ and compute for its surrogate reachset, $\bar{X}\subset \mathbb{R}^{n(\horizon+1)}$. They also compute for another hypercube $\delta X \subset \mathbb{R}^{n(\horizon+1)}$, known as the inflating hypercube, that covers the prediction errors $R^{\mathsf{real},j} , j\in[n\horizon]$, with a provable probabilistic guarantee, and finally propose the following lemma. 
% \begin{lemma}
% Let $\bar{X}$ be a surrogate flowpipe of the surrogate model $\overallf$ for the set of initial conditions $\init$. Let $\PEreal:=\left[ \errreal{1} , \errreal{2}, \cdots , \errreal{n\horizon} \right] $ be the sequence of prediction errors where $\statee_0 \sim \mathcal{W}$, and  let $\delta X$ be the inflating hypercube for $\PEreal$. Then the inflated reachset $X = \bar{X} \oplus \delta X $ is a $\delta$-confident flowpipe for $\trajreal_{\statee_0} \sim \dist$ where $\statee_0 \sim \mathcal{W}$.
% \end{lemma}
% \begin{proof}
%     see \cite{hashemi2024statistical}.
% \end{proof}

Nonetheless, the utilized technique in \cite{hashemi2024statistical} to integrate robust conformal inference, \navid{can be adjusted to reduce the level of conservatism}. Our first contribution in this paper is to address this issue, by combining robust conformal inference with Principal Component Analysis (PCA), and we show this adjustment results in tighter inflating hypercubes and probabilistic reachable sets. \navid{
% This approach also faces scalability challenges when applied to longer-horizon trajectories. This issue arises from their training strategy for the surrogate model. 
In \cite{hashemi2024statistical}, the authors assume a single model that maps the initial state to the entire trajectory\footnote{This approach is motivated by the fact that training a one-step model and iterating it over time leads to the well-known issue of cumulative errors}. However, this results in a large model, which limits the scalability of reachability analysis for extended horizons. Training such a large model becomes inefficient. Furthermore, the model’s large size makes accurate methods for surrogate reachability infeasible, necessitating the use of conservative over approximation techniques.}
% The approaches outlined in \cite{hashemi2024statistical} also encounters scalability challenges when dealing with longer horizon trajectories. This problem stems from their training strategy for the surrogate model. In \cite{hashemi2024statistical}, the authors assume a single model to map the initial state to the entire trajectory\footnote{ To address the problem of cumulative prediction error}. However, this results in a model with large width, limiting scalability for extended horizons due to several factors:
% \begin{itemize}[wide, labelwidth=!, labelindent=0pt]
% \item Training such a model is impractical for longer time horizons or higher system dimensions.
% \item This large width makes exact methods for surrogate reachability infeasible, requiring us to use conservative techniques.
% % \item Using the exact-star method for the surrogate model becomes nearly impossible.
% \end{itemize}
 In response, our second contribution here is to address these challenges by presenting a new training strategy that effectively resolves for all of these listed scalability issues. \navid{We propose training a set of independent small models, each mapping the initial state to a separate sub-partition of the trajectory. This approach eliminates the need for iterating a model over time while maintaining smaller models that are more suitable for efficient training and efficient surrogate reachability analysis.}

\mypara{Related Work}  \navidd{In the context of data-driven reachability analysis, typically, to obtain the reachable set, a dataset of system trajectories needs to be available, e.g., from a black box simulator.} In \cite{devonport2021data}, the authors use Christoffel functions\footnote{See \cite{lasserre2019empirical,marx2021semi} for more details about the Christoffel functions.} to \navidd{learn the model of the system from such a } dataset and generate probabilistic \navidd{reachable} sets. This methodology has been extended in \cite{tebjou2023data} by incorporating conformal inference \cite{vovk2012conditional} \navidd{to improve its data efficiency}. In \cite{devonport2020data}, a Gaussian process-based classifier is employed to learn a model \navidd{that distinguishes} between reachable and unreachable states to approximate the \navidd{reachable} set. \navid{In \cite{devonport2020estimating,dietrich2024nonconvex} \navidd{scenario optimization is used} to generate probabilistic reachable sets.} The method in \cite{fisac2018general} assumes partial knowledge of the \navidd{system} and leverages the dataset to perform statistical reachability analysis. Similarly, the work presented in \cite{dryvr} uses the dataset to \navidd{learn an exponential discrepancy function that estimates trajectory's sensitivity to uncertainty, enabling the computation of probabilistic reachable sets}. Of our particular interest is the method proposed by \cite{ hashemi2024statistical}, which suggests learning a \navidd{system} model on this dataset via $\text{ReLU}$ neural networks \navidd{and then conduct statistical reachability analysis via conformal inference.}  \navidd{By leveraging the ability of neural networks to model complex, high-dimensional relationships alongside the data efficiency of conformal inference, this approach establishes an efficient and structured framework for statistical reachability analysis.} 
% Unlike other approaches, \cite{hashemi2024statistical} focuses on leveraging DNN models to learn from the dataset and also combines this powerful tool with robust conformal inference that is one of the most powerful techniques for statistical error analysis. This powerful combination is the main reason for extending this research rather than focusing on other studies in this paper.
Additionally, the probabilistic guarantees proposed by \cite{hashemi2024statistical} remain valid even when there is a distribution shift between the training and deployment environments. \navidd{This is the main reason we focus on extending this work instead of building on other existing methodologies.}

\navidd{Conformal inference (CI) is a data-efficient method for formally providing guarantees on the $\delta$-quantile of distributions. This method involves sampling an i.i.d. scalar dataset, sorting the samples in ascending order, and demonstrating that one of the sorted samples represents the $\delta$-quantile.} The integration of CI with formal verification techniques has recently received noticeable interest, that is primarily due to its accuracy and level of scalability. For instance, 
% \cite{lin2024verification} combines CI with Hamilton-Jacobi Reachability analysis, a model-based reachability method. 
\cite{bortolussi2019neural} merges CI with neural state classifiers to develop a stochastic runtime verification algorithm. \cite{lindemann2023safe,zecchin2024forking} employ CI to guarantee safety in MPC control using a trained model. \cite{tonkens2023scalable}  applies CI for planning with probabilistic safety guarantees, and \cite{hashemi2024statistical} integrates CI with existing neural network reachability techniques and provides a scalable reachability analysis on stochastic systems, see \cite{lindemann2024formal} for a recent survey article.
 % The authors in \cite{lin2024verification} integrate CI with Hamilton-Jacobi Reachability analysis that is a well-known model based reachability technique. The authors in \cite{bortolussi2019neural} combine CI with Neural state classifiers and introduce a stochastic runtime verification algorithm. The authors in \cite{zecchin2024forking} utilize a trained model for MPC control and rely on CI to propose guarantee on their results, and the authors in \cite{tonkens2023scalable} utilize CI to propose planning with probabilistic safety guarantee.

\mypara{Notation} We use bold letters to represent vectors and vector-valued functions, while \navid{caligraphic} letters  denote sets and distributions. The set $\left\{1,2,\ldots, n \right\}$ is denoted as $[n]$. The Minkowski sum is indicated by $\oplus$. We use $x \sim \mathcal{X}$ to denote that the random variable $x$ is drawn from the distribution $\mathcal{X}$. We present the structure of a feedforward neural network (FFNN) with $\ell$  hidden layers as  an array $[n_0,n_1,\ldots n_{\ell+1}]$, where $n_0$ denotes the number of inputs, $n_{\ell+1}$  is the number of outputs, and $n_i , i\in [\ell]$ denotes the width of the $i$-th hidden layer. We denote $e_i\in\mathbb{R}^n$ as the $i$-th base vector of $\mathbb{R}^n$. We also denote $\lceil x \rceil$ as the smallest integer greater than $x \in \mathbb{R}$.